\section{Разработка слоя клиентского представления}

Интерфейс реализован на серверной стороне с использованием Jinja2 и собственного CSS. Такой подход упрощает сопровождение и позволяет тесно связать формы с логикой бэкенда. Интерфейс ориентирован на студента и обеспечивает быстрый доступ к основным сценариям: вход, регистрация, просмотр расписания, управление дедлайнами.

\subsection{Базовый шаблон и стилизация}

Все страницы наследуются от \texttt{base.html}, который содержит:
\begin{itemize}
    \item верхнюю навигацию (ссылка на календарь, кнопку профиля);
    \item общий CSS-файл \texttt{/static/style.css};
    \item футер с ссылками на API и контактами.
\end{itemize}
В файле \texttt{style.css} описаны стили для карточек задач, модальных окон, аватаров, календаря и компонентов автодополнения.

\subsection{Страницы и шаблоны}

Основные страницы приложения:
\begin{enumerate}
    \item \textbf{Вход} (\texttt{login.html}) --- форма авторизации, отображение ошибок (неверный пароль, неактивный пользователь).
    \item \textbf{Регистрация} (\texttt{register.html}) --- форма создания пользователя с выбором группы и загрузкой аватара. Автодополнение групп реализовано через API \texttt{/api/schedule/groups}.
    \item \textbf{Профиль} (\texttt{profile.html}) --- отображение данных пользователя, редактирование профиля и списка дедлайнов, встроенные модальные окна для создания/редактирования задач.
    \item \textbf{Список задач} (\texttt{tasks\_list.html}) --- карточное отображение задач, быстрые действия (редактировать/удалить).
    \item \textbf{Форма задачи} (\texttt{task\_form.html}) --- создание и редактирование задачи с выбором статуса и дедлайна.
    \item \textbf{Календарь} (\texttt{calendar.html}) --- календарная сетка и расписание на день, выбор пары для автозаполнения дедлайна.
\end{enumerate}

\subsection{Интерактивные компоненты}

Для улучшения UX реализованы следующие элементы:
\begin{itemize}
    \item \textbf{Модальные окна} на страницах профиля и календаря (создание/редактирование задач без перехода на отдельную страницу).
    \item \textbf{Автокомплит предметов и групп} --- заполнение через API расписания, фильтрация по вводу пользователя.
    \item \textbf{Цветовая индикация дедлайнов}: зелёный/жёлтый/красный в зависимости от близости срока.
    \item \textbf{Сортировка задач} на странице профиля (по предмету, срокам, статусу).
\end{itemize}

\subsection{Интеграция календаря с расписанием}

Страница \texttt{calendar.html} загружает полное расписание группы через \texttt{/api/schedule/{group}/full\_schedule}, отображает пары по дням и позволяет создавать задачу кликом на занятие. В форму автоматически подставляются дисциплина и предполагаемый дедлайн по времени пары.
